\documentclass{article}
\usepackage{../../../../doc/latex/cactus}

\newlength{\tableWidth} \newlength{\maxVarWidth} \newlength{\paraWidth} \newlength{\descWidth} \begin{document}

\title{Hydro\_RNSID}
\author{Nik Stergioulas, Roberto De Pietri, Frank L\"ofller}
\date{\today}
\maketitle

\abstract{Hydro\_RNSID - rotating relativistic neutron stars.}

\section{Introduction}
\label{sec:intro}

This thorn generates neutron star initial data for the GRHydro code. As
with the Einstein Toolkit code itself, please feel free to add, alter 
or extend any part of this code. However please keep the documentation up to
date (even, or especially, if it's just to say what doesn't work).

This thorn effectively takes the public domain code {\tt RNSID}
written by Nik Stergioulas and interpolates the output onto a
Cartesian grid. This porting is based on an initila porting to 
Whisky by Luca Baiotti and Ian Hawke and has been adapted to 
GRHydro and Einstein Toolkit.


\section{RNSID}
\label{sec:rnsid}

RNSID, or rotating neutron star initial data, is a code based on the
Komatsu-Eriguch-Hachisu (KEH) method for constructing models of
rotating neutron stars. It allows for polytropic or tabulated
equations of state. For more details of the how the code works
see~\cite{rnsid1}, \cite{nsthesis} (appendix A is particularly
helpful) or especially \cite{nslivrev} which is the most up to date
and lists other possible methods of constructing rotating neutron star
initial data.

In short Hydro\_RNSID is a thorn that generate initial model for rotating 
isolated stars described by a zero-temperature tabulated Equation of State or
an iso-entrophic politropic EOS. The activation of the thorn for
genereting ID (The thorns ``Hydro\_Base'' and  ``GRHydro'' are the two prerequisites)

The model are generated specifing the central baryonic density
(rho\_central), the oblatness of the Star (axes\_ratio) and the
rotational profile (rotation\_type).  Currently two kinds of
rotational profiles are implemented: ``uniform'' for uniformly
rotating stars and ``diff'' for differentially rotating stars,
described by the j-law profile (parametrized by the parameter
A\_diff=$\hat{A}$):
\begin{equation}
  \Omega_c-\Omega = \frac{1}{\hat{A}^2 r_e^2}
  \left[ \frac{(\Omega-\omega) r^2 \sin^2 \theta e^{-2\nu}}{1-(\Omega-\omega)^2 r^2 \sin^2 \theta e^{-2\nu}}\right]
\end{equation}
where $r_e$ is the equatorial radius of the star and $\Omega$ is the
rotational angular velocity $\Omega=u^\phi/u^0$ and $\Omega_c$ is
$\Omega$ at the center of the star.


\section{Parameters of Thorn}
\label{sec:rnsid_par}

Here one can find definition of the main parameter the determine the
behaviour of the Thorn. The activation of the RNSID initial data
is achieved by the following line:
\begin{verbatim}
ActiveThorns="Hydro_Base GRHydro Hydro_RNSID"
#####
##### Setting for activating the ID 
#####
ADMBase::initial_data  = "hydro_rnsid"
ADMBase::initial_lapse = "hydro_rnsid"
ADMBase::initial_shift = "hydro_rnsid"
\end{verbatim}

The correspongig section of the parameter file is:
\begin{verbatim}
#####
##### Basic Setting
#####
Hydro_rnsid::rho_central   = 1.28e-3  # central baryon density (G=c=1) 
Hydro_rnsid::axes_ratio    = 1        # radial/equatorial axes ratio 
Hydro_rnsid::rotation_type = diff     # uniform = uniform rotation
Hydro_rnsid::A_diff        = 1        # Parameter of the diff rot-law.
Hydro_rnsid::accuracy      = 1e-10    # accuracy goal for convergence
\end{verbatim}



Than a section for setting the Equation of State (EOS) should be added.
If this section is missing a ``poly'' EOS will be used with default parameters.
The two possibilities are:
\begin{description}
\item{\textbf{Isentropic Polytrope:}}
  
  In this case the base setting for the initial data are specified
  giving the following parameters:
\begin{verbatim}
#####
##### Setting for polytrope
#####
Hydro_rnsid::eos_type  = "poly"  
Hydro_rnsid::RNS_Gamma = 2.0
Hydro_rnsid::RNS_K     = 165   
\end{verbatim}
They correspond at the following implementation of the EOS
that it is consistent with the 1st Law of thermodinamics.
\begin{eqnarray}
   p         &=& K \cdot \rho^\Gamma \\
   \epsilon &=& \frac{K}{\Gamma-1} \cdot \rho^{\Gamma-1} 
\end{eqnarray}
and for the above choice of parameters corresponding to the choice $K=165$ (in units where $G=c=M_\odot=1$) and $\Gamma=2$.

\item{\textbf{Tabulated EOS:}}

  In this case the (cold) EOS used to
  generate the initial data is read from a file

\begin{verbatim}
#####
##### Setting for tabulated EOS
#####
Hydro_rnsid::eos_type  = "tab"   
Hydro_rnsid::eos_file  = "full_path_name_of_the tabulated_EOS_file"
\end{verbatim}
  The syntax of the tabulated file is the same as for the original
  RNSID program and assumes that all quantities are expressed in the
  cgs system of units. The first line contains the number of tabulated
  values ($N$) while the next $N$ lines contain the values:
  $e=\rho(1+\varepsilon)$, $p$, $\log h = c^2 \log_e((e+p)/\rho)$, and
  $\rho$, respectively.
\end{description}


An additional section allows one to start initial data from a
previously generated binary file or to save the data generated at this
time. Usually the best way to proceed is to specify where the initial
data file should be located.
\begin{verbatim}
#####
##### Setting for recover and saving of 2d models 
#####
Hydro_rnsid::save_2Dmodel    = "yes"   # other possibility is no (default)
Hydro_rnsid::recover_2Dmodel = "yes"   # other possibility is no (default)
Hydro_rnsid::model2D_file    = "full_file_name"
\end{verbatim}

For examples of initial data generated using RNSID and their
evolutions, see~\cite{rotns1,rotns2}. In the par directory, examples
are provided as a perl file that produces the corresponding Cactus par
files. These examples correspond to the evolutions described
in~\cite{rotns1,rotns2}.

\section{Utility program}

Together with the Thorn, we distribute a self-executable version of
the initial data routine RNSID that accepts the same parameters as the
thorn and is able to create a binary file of the 2d initial data that
can be directly imported into the evolution code. Moreover, the
program RNS\_readID is provided that reads a 2d initial data file and
produces an hdf5 version of the data interpolated onto a 3d grid.


\begin{thebibliography}{1}

\bibitem{rotns1}
J.~A. Font, N. Stergioulas and K.~D. Kokkotas.
\newblock Nonlinear hydrodynamical evolution of rotating relativistic
stars: Numerical methods and code tests.
\newblock {\em Mon. Not. Roy. Astron. Soc.}, {\bf 313}, 678, 2000.

\bibitem{rotns2}
F.~L\"offler, R.~De Pietri, A.~Feo, F.~Maione and L.~Franci, 
\newblock Stiffness effects on the dynamics of the bar-mode instability 
of neutron stars in full general relativity. 
\newblock {\em Phys. Rev.}, {\bf D 91}, 064057, 2015 (arXiv:1411.1963).

\bibitem{rnsid1}
N. Stergioulas and J.~L. Friedmann.
\newblock Comparing models of rapidly rotating relativistic stars
constructed by two numerical methods.
\newblock {\em ApJ.}, {\bf 444}, 306, 1995.

\bibitem{nsthesis}
N. Stergioulas.
\newblock The structure and stability of rotating relativistic stars.
\newblock PhD thesis, University of Wisconsin-Milwaukee, 1996.

\bibitem{nslivrev}
N. Stergioulas.
\newblock Rotating Stars in Relativity
\newblock {\em Living Rev. Relativity}, {\bf 1}, 1998.
\newblock [Article in online journal], cited on 18/3/02,
  http://www.livingreviews.org/Articles/Volume1/1998-8stergio/index.html.

\end{thebibliography}


\include{interface}
\include{param}
\include{schedule}


\end{document}
